\documentclass[12pt,letterpaper]{article}

% ---- Packages ----
\usepackage[utf8]{inputenc}
\usepackage[T1]{fontenc}
\usepackage{amsmath,amssymb}
\usepackage{newtxtext,newtxmath}
\usepackage[margin=1in]{geometry}
\usepackage{setspace}
\usepackage{titlesec}
\usepackage{enumitem}
\usepackage{hyperref}
\usepackage{microtype}
\usepackage{booktabs}

% ---- Section formatting ----
\titleformat{\section}{\large\bfseries\sffamily}{\thesection.}{0.5em}{}
\titleformat{\subsection}{\normalsize\bfseries\sffamily}{\thesubsection}{0.5em}{}
\titleformat{\subsubsection}{\normalsize\bfseries}{\thesubsubsection}{0.5em}{}

% ---- Spacing ----
\setstretch{1.15}
\setlength{\parskip}{0.5em}
\setlength{\parindent}{0em}

% ---- Custom commands ----
\newcommand{\rhoa}{\rho_{\!A}}
\newcommand{\Ka}{K_{\!A}}
\newcommand{\rloop}{r_{\mathrm{loop}}}

\begin{document}

\begin{center}
{\LARGE\bfseries Unifying Theory of Dynamic Aether Mechanics:\\
Bound Core Electron Model}\\[1em]
{\large Erik Otto}\\[0.5em]
{\normalsize\itshape Working Paper in the Dynamic Aether Mechanics Framework}
\end{center}

\vspace{1em}

\begin{abstract}
Paper~7 \cite{otto-quantum} identifies the electron as a quantized vortex ring
in the Aether condensate, with mass determined by the Donnelly kinetic energy
formula.\ The vortex core is treated as a depleted region where the condensate
density drops to zero at the healing length~$\xi$.\ This paper proposes a
refinement: the Dynamic Casimir Effect (DCE), which operates at the vortex
boundary, creates bound Aether that migrates inward under the Bernoulli
pressure gradient and accumulates as a toroidal ring in the vortex core.\ The
bound ring grows until the boundary velocity drops below the DCE permanence
threshold, at which point only transient binding (gravity) occurs and the
electron mass stabilizes.\ This self-regulating mechanism explains why the
electron has the mass it does.\ The self-consistent vortex equation, corrected
for core energy ($+5$--$14\%$) and the Bernoulli density deficit ($-1.7\%$),
predicts $m_e \approx 0.4$--$0.6$~MeV across the allowed parameter range and
matches the observed $0.511$~MeV for $m_A \approx 616$--$644$~MeV/$c^2$, well
within the independently constrained convergence zone.\ A companion
derivation obtains Coulomb's constant from the ring's spin-phase winding
pump ($k e^2 = \rhoa \kappa_e^2 \rloop^2/8\pi$) and the Bohr magneton with
$g = 2$ exactly from the same geometry, tying the electric, magnetic, and
mass sectors to a single topological postulate: one winding advance per
Compton cycle.
\end{abstract}

\vspace{1em}

\bigskip
\noindent\textit{This working paper explores a specific mechanism within the
Dynamic Aether Mechanics framework.\ It references results from Paper~1
(General Overview)~\cite{otto-overview}, Paper~2 (Special
Relativity)~\cite{otto-relativity}, Paper~3
(Gravity)~\cite{otto-gravity}, Paper~4 (Magnetic and Electric
Fields)~\cite{otto-em}, Paper~5 (Strong and Weak
Forces)~\cite{otto-strong}, and Paper~7 (Quantum
Mechanics)~\cite{otto-quantum} of the main series.\ The ideas presented here
are under development and may be incorporated into the main papers in revised
form.}
\bigskip


% ============================================================
\section{Introduction}
\label{sec:introduction}
% ============================================================

Paper~7 \cite{otto-quantum} establishes that the electron is a quantized
vortex ring in the Aether condensate with circulation
$\kappa_e = h/(2m_A)$.\ The electron rest energy is set equal to the Donnelly
kinetic energy of the circulatory flow:
%
\begin{equation}
  m_e\,c^2 = \tfrac{1}{2}\,\rhoa\,\kappa_e^2\,\rloop
  \left[\ln\!\left(\frac{8\,\rloop}{\xi}\right) - \tfrac{1}{2}\right]
  \label{eq:paper7-mass}
\end{equation}
%
where $\rloop = \hbar/(m_e c)$ is the loop radius (reduced Compton
wavelength), $\xi = \hbar/(m_A c\sqrt{2})$ is the healing length, and the
core is treated as a depleted region where the condensate density drops to
zero.

This treatment leaves two questions open:

\begin{enumerate}[leftmargin=2em]
  \item \textbf{Where does the DCE product go?}\ Paper~1
    \cite{otto-overview} establishes that the spinning vortex core drives the
    Dynamic Casimir Effect: transient binding produces gravity, and permanent
    binding (above the DCE threshold) produces mass.\ But the permanently
    bound material must go somewhere.\ The current model is silent on its
    fate.

  \item \textbf{Why does the electron have the mass it does?}\
    Eq.~(\ref{eq:paper7-mass}) determines $m_e$ given $m_A$ and $\rhoa$, but
    provides no mechanism for \emph{why} this particular vortex energy is
    selected from among all possible configurations.
\end{enumerate}

This paper proposes that both questions have the same answer: the DCE
permanently binds Aether at the vortex boundary, the bound material migrates
inward under the Bernoulli pressure gradient and accumulates as a toroidal
ring in the core, the ring grows until the boundary velocity drops below the
permanence threshold, and the electron mass stabilizes at this
self-regulating equilibrium.


% ============================================================
\section{Formation Model}
\label{sec:formation}
% ============================================================

The electron forms through the following sequence:

\begin{enumerate}
  \item \textbf{Vortex ring creation.}\ In the high-energy early universe,
    turbulent Aether flows create quantized vortex rings with the minimum
    (spin-$\tfrac{1}{2}$) circulation quantum
    $\kappa_e = h/(2m_A)$.\ Initially these are pure fluid vortices with no
    bound material---the core is depleted at the healing length as in the
    standard model.

  \item \textbf{DCE creates matter.}\ The circulation velocity at the
    healing length is (Section~\ref{sec:circulation-velocity}):
    %
    \begin{equation}
      v(\xi) = \frac{c}{\sqrt{2}} \approx 0.707\,c
      \label{eq:v-at-xi}
    \end{equation}
    %
    This exceeds the DCE permanence threshold, so the rapidly-circulating
    boundary continuously converts vacuum fluctuations into permanent bound
    Aether \cite{wilson2011}.

  \item \textbf{Matter migrates inward.}\ The Bernoulli pressure in the
    vortex is lowest at the core
    (Section~\ref{sec:bernoulli-pressure}).\ Bound material, no longer part
    of the superfluid flow, is pushed inward by the pressure gradient and
    accumulates along the vortex tube axis, forming a toroidal ring of bound
    Aether tracing the major radius $\rloop$.

  \item \textbf{Ring grows.}\ As more bound material falls to the center,
    the ring expands radially outward from the tube axis.\ The effective core
    boundary moves from $\xi$ to a larger radius $r_{\mathrm{eq}}$.

  \item \textbf{Boundary velocity drops.}\ The circulation $\kappa_e$ is
    topologically quantized and fixed, but the velocity at the boundary is
    $v(r) = \kappa_e/(2\pi r)$.\ As $r_{\mathrm{eq}}$ increases, the
    boundary velocity decreases.

  \item \textbf{Permanence threshold reached.}\ When
    $v(r_{\mathrm{eq}})$ drops to $v_{\mathrm{perm}}$---the minimum velocity
    for permanent DCE binding---the vortex can no longer create permanent
    mass.\ Only transient bind--unbind cycling occurs at the boundary,
    producing the asymmetric pressure oscillation that manifests as gravity
    \cite{otto-gravity}.

  \item \textbf{Mass stabilizes.}\ The electron's mass is the sum of the
    remaining circulation kinetic energy and the bound ring's rest mass at
    this equilibrium point.
\end{enumerate}

This is a \emph{self-regulating} process: the creation of mass enlarges the
core, which reduces the boundary velocity, which eventually shuts off mass
creation.\ The electron mass is not a free parameter but an equilibrium
outcome determined by medium properties.

\subsection{Energy Conservation}
\label{sec:energy-conservation}

The DCE mechanism converts energy from the vortex's circulation into bound
material---the moving boundary does work on the vacuum fluctuations, and
the energy is drawn from the kinetic energy of the flow.\ As the ring grows,
the circulation KE decreases by the same amount that the ring mass increases.\
The total energy is approximately conserved during the formation process:
%
\begin{equation}
  m_e\,c^2 \;\approx\; E_{\mathrm{KE,initial}}
  \;=\; E_{\mathrm{KE,final}} + m_{\mathrm{ring}}\,c^2
  \label{eq:energy-conservation}
\end{equation}
%
The electron's final mass is set by the initial vortex energy at creation,
not by the partitioning between KE and ring mass.\ The circulation quantum
$\kappa_e = h/(2m_A)$ is topologically fixed and cannot change; the KE
reduction comes from the shrinking integration domain as the core expands
from $\xi$ to $r_{\mathrm{eq}}$.


% ============================================================
\section{Key Equations}
\label{sec:equations}
% ============================================================

\subsection{Circulation Velocity}
\label{sec:circulation-velocity}

For a quantized vortex with circulation
$\kappa_e = h/(2m_A) = \pi\hbar/m_A$, the flow velocity at distance $r$
from the tube axis is:
%
\begin{equation}
  v(r) = \frac{\kappa_e}{2\pi\,r}
  \label{eq:v-of-r}
\end{equation}

At the healing length $\xi = \hbar/(m_A c\sqrt{2})$, the velocity evaluates
to $v(\xi) = c/\sqrt{2} \approx 0.707\,c$ (Eq.~\ref{eq:v-at-xi}).\ This
result is exact and independent of $m_A$: the circulation quantum and the
healing length are both proportional to $1/m_A$, so $m_A$ cancels.

\subsection{Bernoulli Pressure}
\label{sec:bernoulli-pressure}

The pressure distribution in the vortex flow:
%
\begin{equation}
  P(r) = P_\infty - \tfrac{1}{2}\,\rhoa\,v^2(r)
       = P_\infty - \frac{\rhoa\,\kappa_e^2}{8\pi^2\,r^2}
  \label{eq:bernoulli}
\end{equation}
%
Pressure is lowest at the core center and increases outward.\ Bound material
(not participating in superfluid flow) experiences a net inward force from
the surrounding higher-pressure fluid, driving inward migration.

\subsection{Equilibrium Condition}
\label{sec:equilibrium}

The bound ring grows until the boundary velocity reaches the permanence
threshold:
%
\begin{equation}
  v(r_{\mathrm{eq}}) = v_{\mathrm{perm}}
  \quad\Longrightarrow\quad
  r_{\mathrm{eq}} = \frac{\kappa_e}{2\pi\,v_{\mathrm{perm}}}
  \label{eq:req}
\end{equation}
%
For $v_{\mathrm{perm}} < c/\sqrt{2}$, the ring extends beyond the healing
length: $r_{\mathrm{eq}} > \xi$.



% ============================================================
\section{The Electron Mass Equation}
\label{sec:mass-equation}
% ============================================================

\subsection{Core Energy}
\label{sec:core-energy}

The Donnelly formula \cite{donnelly1991} (Eq.~\ref{eq:paper7-mass}) computes
only the kinetic energy of the circulating flow outside the core.\ In a standard GPE vortex,
the total energy also includes the core's gradient and potential energy---the
cost of depleting the condensate.\ This core energy is of order
$1/\ln(\rloop/\xi)$ of the kinetic energy, or $\sim\!5$--$14\%$ for the
electron's aspect ratio (the exact value for the logarithmic potential used
in this theory has not been computed from first principles).

In the bound-core model, the core region is filled with bound material
instead of depleted vacuum.\ The bound ring's rest mass \emph{replaces} the
gradient and potential energy of the depleted core---it is the physical
realization of the same core energy, not an additional contribution on top
of it.

\subsection{Bernoulli Density Deficit}
\label{sec:bernoulli-density}

The Donnelly formula assumes constant density $\rhoa$ everywhere.\ But the
vortex circulation itself produces a Bernoulli density deficit.\ This is a
genuine density change (not just a pressure change with constant $\rhoa$)
because the vortex circulation is \emph{driven} flow, not free-fall---the
persistent spin energy maintains the flow against Bernoulli pressure,
creating a venturi-like deficit.\ This is physically distinct from
gravitational inflow, where free-fall dynamics ensure $\Delta\rho = 0$
exactly \cite{otto-relativity}:
%
\begin{equation}
  \frac{\delta\rho(r)}{\rhoa} \approx -\frac{v^2(r)}{2\,c^2}
  = -\frac{\kappa_e^2}{8\pi^2\,c^2\,r^2}
  \label{eq:density-deficit}
\end{equation}

At the healing length ($v = c/\sqrt{2}$), the deficit reaches
$\delta\rho/\rhoa \approx -25\%$, dropping rapidly with distance:

\begin{center}
\begin{tabular}{@{}ccc@{}}
\toprule
$r/\xi$ & $v/c$ & $\delta\rho/\rhoa$ \\
\midrule
1.0 & 0.707 & $-25\%$ \\
1.5 & 0.471 & $-11\%$ \\
2.0 & 0.354 & $-6\%$ \\
5.0 & 0.141 & $-1\%$ \\
\bottomrule
\end{tabular}
\end{center}

The density deficit reduces the kinetic energy.\ Integrating the correction
$\delta E = \int\!\tfrac{1}{2}\,(-\delta\rho)\,v^2\,dV$:
%
\begin{equation}
  \frac{\delta E_{\mathrm{KE}}}{E_{\mathrm{KE}}} \approx
  -\frac{1}{8\,\ln(\rloop/\xi)}
  \approx -1.7\%
  \label{eq:ke-correction}
\end{equation}
%
The factor $1/8$ arises because $\kappa_e^2/(16\pi^2 c^2 \xi^2) = 1/8$ for
$v(\xi) = c/\sqrt{2}$.\ For the logarithmic equation of state, the exact
compressible Bernoulli gives $\rho(r) = \rhoa\,e^{-v^2/(2c^2)}$ rather than
the linearized $\rhoa\,[1 - v^2/(2c^2)]$; the integrated KE correction
differs by less than $3\%$ between exact and linear forms, so the linearized
result is adequate.

The Donnelly formula itself is non-relativistic.\ In the Aether framework,
relativistic effects emerge from the medium at $O(v^2/c^2)$: time dilation
from wave path geometry, length contraction from density changes
\cite{otto-relativity}.\ The Bernoulli density correction is the leading
such correction to the vortex energy; higher-order terms are
$O(v^4/c^4) \sim 0.4\%$, negligible at this level of precision.

\subsection{Combined Mass Equation}
\label{sec:combined}

Combining all contributions:
%
\begin{equation}
  \boxed{m_e\,c^2 \;\approx\;
  \underbrace{\tfrac{1}{2}\,\rhoa\,\kappa_e^2\,\rloop
  \left[\ln\!\left(\frac{8\,\rloop}{r_{\mathrm{eq}}}\right)
  - \tfrac{1}{2}\right]}_{\displaystyle E_{\mathrm{KE}}\;\text{(circulation)}}
  \;\times\;\underbrace{(1 - 0.017)}_{\text{Bernoulli}}
  \;+\; \underbrace{E_{\mathrm{core}}}_{\displaystyle\sim\!5\text{--}14\%
  \text{ of }E_{\mathrm{KE}}}}
  \label{eq:me-total}
\end{equation}

The net correction from the bare Donnelly formula is:
%
\begin{equation}
  m_e\,c^2 \approx E_{\mathrm{KE}}^{\mathrm{Donnelly}} \times
  (1 - 0.017 + f_{\mathrm{core}})
  \label{eq:net-correction}
\end{equation}
%
where $f_{\mathrm{core}} \approx 0.05$--$0.14$ is the core energy fraction
(model-dependent).\ The net correction ranges from $+3\%$ to $+12\%$
above the bare Donnelly value.\ The loop radius is
$\rloop = \hbar/(m_e\,c)$ (the reduced Compton wavelength), following
Paper~7~\cite{otto-quantum}.


% ============================================================
\section{Mass Predictions}
\label{sec:predictions}
% ============================================================

\subsection{Donnelly Kinetic Energy}
\label{sec:ke-only}

Using the self-consistent condition $\rloop = \hbar/(m_e c)$ and setting
$r_{\mathrm{eq}} = \xi$, the Donnelly formula gives:
%
\begin{equation}
  m_e^2 = \frac{\pi^2\,\rhoa\,\hbar^3}{2\,m_A^2\,c^3}
  \left[\ln\!\left(\frac{8\sqrt{2}\,m_A}{m_e}\right) - \tfrac{1}{2}\right]
  \label{eq:self-consistent}
\end{equation}

This transcendental equation is solved iteratively.\ Results for
$\rhoa = 5 \times 10^{11}$~kg/m$^3$:

\begin{center}
\begin{tabular}{@{}ccccc@{}}
\toprule
$m_A$ & $\xi$ & $\rloop$ & $m_e$ (KE only) & vs.\ observed \\
(MeV/$c^2$) & (fm) & (fm) & (MeV) & \\
\midrule
500 & 0.28 & 325 & 0.61 & $+19\%$ \\
550 & 0.25 & 354 & 0.56 & $+9\%$ \\
600 & 0.23 & 383 & 0.515 & $+1\%$ \\
650 & 0.21 & 411 & 0.48 & $-6\%$ \\
700 & 0.20 & 439 & 0.45 & $-12\%$ \\
800 & 0.17 & 495 & 0.40 & $-22\%$ \\
\bottomrule
\end{tabular}
\end{center}

The bare Donnelly formula gives $m_e \approx 0.4$--$0.6$~MeV across the
full convergence zone ($500$--$850$~MeV/$c^2$ from
Paper~7~\cite{otto-quantum}), bracketing the observed $0.511$~MeV.

\subsection{Corrected Prediction}
\label{sec:corrected}

Since DCE converts circulation KE into ring mass (energy conserved,
Eq.~\ref{eq:energy-conservation}), the total electron mass is set by the
initial vortex energy.\ The Donnelly formula gives only the KE portion.\
Including the core energy and Bernoulli corrections
(Eq.~\ref{eq:net-correction}), the corrected prediction depends on the core
energy fraction $f_{\mathrm{core}}$.

One propagation subtlety must be handled carefully.\ The energy relation
$m_e c^2 = E_{\mathrm{KE}} \times (1 + \mathrm{net})$ does not shift the
self-consistent mass linearly, because
$E_{\mathrm{KE}} \propto \rloop \propto 1/m_e$: substituting
$\rloop = \hbar/(m_e c)$ places the correction in the \emph{squared}
mass, $m_e^2 \propto (1 + \mathrm{net})$, so the self-consistent solution
shifts by $\sqrt{1 + \mathrm{net}}$---half the naive exponent:

\begin{center}
\begin{tabular}{@{}ccccc@{}}
\toprule
$f_{\mathrm{core}}$ & Net correction & Best-fit $m_A$ & KE only (MeV) &
Corrected (MeV) \\
\midrule
$5\%$ & $+3\%$ & 616~MeV & 0.503 & 0.511 \\
$8\%$ & $+6\%$ & 625~MeV & 0.496 & 0.511 \\
$11\%$ & $+9\%$ & 635~MeV & 0.490 & 0.511 \\
$14\%$ & $+12\%$ & 644~MeV & 0.484 & 0.511 \\
\bottomrule
\end{tabular}
\end{center}

For any core energy in the $5$--$14\%$ range, there exists an $m_A$ in the
convergence zone that matches the observed mass:
%
\begin{equation}
  \boxed{m_e \approx 0.511\;\text{MeV} \quad\text{(predicted, }m_A \approx
  616\text{--}644\;\text{MeV)}
  \qquad\text{vs.}\qquad
  m_e = 0.511\;\text{MeV} \quad\text{(observed)}}
  \label{eq:result}
\end{equation}

The Bernoulli density deficit ($-1.7\%$) partially offsets the core energy
($+5$--$14\%$), and the remaining net correction
($+3$--$12\%$)---entering as $\sqrt{1+\mathrm{net}}$---is absorbed by
$m_A$, which shifts from $600$~MeV (bare Donnelly) to $616$--$644$~MeV
(corrected).


% ============================================================
\section{Physical Structure at Equilibrium}
\label{sec:structure}
% ============================================================

\subsection{Compression Ratio}
\label{sec:compression}

Even if the entire electron mass were packed into the core ring (the extreme
$f = 1$ case), the compression ratio would be remarkably low:
%
\begin{equation}
  X = \frac{\rho_{\mathrm{core}}}{\rhoa}
    = \frac{m_e}{\rhoa\,\cdot\,2\pi^2\,\rloop\,\xi^2}
    \approx 4\text{--}5
  \label{eq:compression}
\end{equation}
%
(varying from $\sim\!4.4$ at $m_A = 600$ to $\sim\!4.9$ at $m_A = 630$).\
Only $\sim\!5$ times ambient density even in the extreme case---far below the
$\sim\!7 \times 10^7$ compression in QCD bridges \cite{otto-strong}.\ This
reflects the electron's small mass ($0.511$~MeV) spread over a large
circumference ($2\pi\rloop \approx 2400$~fm).

For the core energy range $f_{\mathrm{core}} = 5$--$14\%$, the ring carries
only $0.03$--$0.07$~MeV.\ Spread over the core volume
($\sim\!300$--$400$~fm$^3$), the ring density is well below $\rhoa$---the
core is sparsely filled, not compressed.

\subsection{Self-Regulation}
\label{sec:self-regulation}

The self-regulation mechanism proceeds as follows:

\begin{enumerate}[leftmargin=2em]
  \item DCE creates bound material at the core boundary where
    $v > v_{\mathrm{perm}}$
  \item Material falls inward, expanding the effective core from $\xi$ to
    $r_{\mathrm{eq}}$
  \item At $r_{\mathrm{eq}}$, the velocity has dropped:
    $v(r_{\mathrm{eq}}) < v(\xi)$
  \item When $v(r_{\mathrm{eq}}) = v_{\mathrm{perm}}$, permanent binding
    ceases
  \item Transient binding continues (producing gravity), but no new mass
    accumulates
\end{enumerate}

The quantized circulation $\kappa_e = h/(2m_A)$ is topologically fixed---it
cannot change.\ Only the boundary radius changes as the ring grows, reducing
the velocity through the $v \propto 1/r$ dependence.

Under energy conservation (Eq.~\ref{eq:energy-conservation}), the total
electron mass is fixed by the initial vortex energy at creation.\ The
threshold velocity $v_{\mathrm{perm}}$ determines only the partitioning
between circulation KE and ring rest mass, not the total.\ The ring mass at
equilibrium is:
%
\begin{equation}
  m_{\mathrm{ring}}\,c^2 = f_{\mathrm{core}} \times m_e\,c^2
  \label{eq:mring}
\end{equation}
%
where $f_{\mathrm{core}} \approx 5$--$14\%$.\ The electron remains
circulation-dominated at equilibrium.

\subsection{The DCE Permanence Threshold}
\label{sec:threshold}

The permanence threshold $v_{\mathrm{perm}}$ determines the equilibrium core
size $r_{\mathrm{eq}} = \kappa_e/(2\pi v_{\mathrm{perm}})$, the KE/ring
mass partitioning, and the detailed structure of the electron.\ For
$v_{\mathrm{perm}}$ near $c/\sqrt{2}$ (modest ring growth), the total mass
varies only weakly with $v_{\mathrm{perm}}$.\ A natural estimate from the
Schwinger pair-production analogy gives:
%
\begin{equation}
  v_{\mathrm{perm}} \sim \frac{c}{\sqrt{\pi}} \approx 0.56\,c
  \label{eq:threshold-estimate}
\end{equation}
%
This corresponds to $r_{\mathrm{eq}} \approx 1.26\,\xi$---a modest expansion
of the core beyond the healing length.\ Quantifying $v_{\mathrm{perm}}$ from
first principles is the key open calculation
(Section~\ref{sec:open-questions}).


% ============================================================
\section{Significance}
\label{sec:significance}
% ============================================================

\subsection{Electron Mass from Medium Properties}

The electron mass is determined by $m_A$, $\rhoa$, and $c$---all
properties of the Aether medium---through a specific transcendental equation
(Eq.~\ref{eq:self-consistent}) with no adjustable structure.\ The mass
arises from three interlocking constraints:

\begin{itemize}[leftmargin=2em]
  \item \textbf{Circulation quantization}: $\kappa_e = h/(2m_A)$ fixes the
    flow speed at each radius
  \item \textbf{Quantum minimum size}: $\rloop = \hbar/(m_e c)$ identifies
    the loop radius with the reduced Compton wavelength
    (Paper~7~\cite{otto-quantum})
  \item \textbf{Self-consistency}: $m_e$ must equal the energy of a vortex
    ring of size $\rloop$, which itself depends on $m_e$
\end{itemize}

For $m_A \approx 616$--$644$~MeV (within the convergence zone), the
corrected prediction matches the observed $0.511$~MeV.\ The inputs $m_A$ and
$\rhoa$ are independently constrained by QCD bridge physics
\cite{otto-strong} and the healing-length convergence
\cite{otto-quantum}.\ The core energy fraction $f_{\mathrm{core}}$ is the
remaining model-dependent quantity; computing the vortex core structure in the
logarithmic potential would pin down both $f_{\mathrm{core}}$ and the
best-fit $m_A$ simultaneously.

\subsection{Self-Regulating Mass}

The bound-core model explains \emph{why} the electron has the mass it
does: the DCE mechanism creates mass until it shuts itself off.\ Every
electron reaches the same equilibrium because every minimum-circulation
vortex ring has the same $\kappa_e$, the same velocity profile, and
therefore the same threshold crossing point.\ Mass universality follows from
circulation quantization.

\subsection{Open Tube / Closed Tube Taxonomy}

The bound electron core ring is a \emph{closed tube} of gently compressed
Aether, stabilized by the encircling vortex circulation.\ The QCD bridge is
an \emph{open tube} of strongly compressed Aether, stabilized by quark
vortex endpoints \cite{otto-strong}.\ Both are topologically-stabilized bound
Aether structures in the same medium, differing in topology (closed vs.\
open), compression ($\ll 1$ to $\sim\!5$ vs.\ $\sim\!7 \times 10^7$), and
energy scale ($0.5$~MeV vs.\ $\sim\!1$~GeV/fm).

\subsection{What Is Derived vs.\ Assumed}

The self-consistent mass equation (Eq.~\ref{eq:self-consistent}) is derived
from the Donnelly vortex energy formula \cite{donnelly1991}, circulation
quantization ($\kappa_e = h/(2m_A)$), and the identification
$\rloop = \hbar/(m_e c)$ from Paper~7~\cite{otto-quantum}.\ The Bernoulli
correction (Eq.~\ref{eq:ke-correction}) is derived from the
density-pressure coupling in the unbound superfluid.\ The formation narrative
and energy conservation (Eq.~\ref{eq:energy-conservation}) are derived from
the DCE mechanism and Bernoulli pressure gradient.

Three quantities remain undetermined from first principles: $m_A$
(constrained to $500$--$850$~MeV/$c^2$ by QCD bridge physics and
healing-length convergence, but not uniquely determined),
$f_{\mathrm{core}}$ (the core energy fraction, which requires solving the
GPE with the logarithmic potential), and $\rloop = \hbar/(m_e c)$ (the loop
radius, identified with the Compton wavelength in Paper~7 but not derived
from vortex structure).\ Deriving $\rloop$ from first principles remains an
open problem (Section~\ref{sec:open-questions}).\ Independently determining
$m_A$ or $f_{\mathrm{core}}$ would reduce the remaining freedom.


% ============================================================
\section{Relationship to Paper~7}
\label{sec:relationship}
% ============================================================

This working paper proposes three modifications to the electron model in
Paper~7:

\begin{enumerate}[leftmargin=2em]
  \item \textbf{Core interpretation.}\ The vortex core is filled with bound
    Aether rather than being a depleted vacuum.\ The condensate density still
    drops to zero at $\sim\!\xi$; the core region contains non-condensate
    material (bound Aether) rather than nothing.

  \item \textbf{Energy corrections.}\ The Donnelly formula undercounts the
    total energy by $\sim\!5$--$14\%$ (core energy, model-dependent) and
    overcounts by $\sim\!1.7\%$ (Bernoulli density deficit).\ The net
    correction, entering the self-consistent mass as
    $\sqrt{1+\mathrm{net}}$, shifts the best-fit $m_A$ from
    $\sim\!600$~MeV to $\sim\!616$--$644$~MeV, still within the
    convergence zone.

  \item \textbf{Mass origin.}\ The electron mass is not merely
    \emph{determined by} the vortex equation---it is \emph{selected} by the
    self-regulating DCE mechanism.\ The formation narrative provides a
    physical reason why the vortex settles at this particular energy.

\end{enumerate}

No equations in Paper~7 need to be changed.\ The Donnelly formula remains
the dominant contribution ($\sim\!85$--$95\%$ of the total mass).\ The
corrections refine the numerical prediction and shift the optimal $m_A$
within the existing allowed range.


% ============================================================
\section{The Form Factor Challenge}
\label{sec:form-factor}
% ============================================================

Electron scattering experiments (Bhabha scattering at LEP, M\o ller
scattering, $g\!-\!2$ measurements) find no deviation from the QED
point-particle prediction up to momentum transfers
$Q \sim 200$~GeV, corresponding to a spatial resolution of
$\sim\!10^{-18}$~m \cite{LEP-Bhabha}.\ Yet the vortex model predicts an
electron with loop radius $\rloop \approx 386$~fm
$= 3.9 \times 10^{-13}$~m---five orders of magnitude larger.\ This section
examines what the experiments actually constrain and how the vortex model
may be consistent.

\subsection{What the Experiments Measure}

These experiments do not image the electron.\ They measure scattering
cross-sections and compare them to QED predictions for a structureless
point particle.\ ``Point-like to $10^{-18}$~m'' means: \emph{the
electron's electromagnetic interactions match QED predictions at all tested
momentum transfers.}\ Any internal structure would modify the
electron--photon vertex, producing a form factor $F(Q^2) \neq 1$ that
changes the cross-section at the scale $Q \sim \hbar c / r$, where $r$
is the charge radius.

For comparison, the proton has a charge radius of $0.87$~fm and its form
factor deviates from~$1$ at $Q \sim 200$~MeV, exactly as expected from
its spatial extent.\ The proton's charge radius ($0.87$~fm) exceeds its
Compton wavelength ($0.21$~fm) by a factor of~$4$, so structure is always
resolvable.

For the electron vortex, if the charge were distributed around the
$386$~fm ring, $F(Q^2)$ would deviate from~$1$ at $Q \sim 0.5$~MeV.\
This is six orders of magnitude below LEP's reach---the deviation would be
enormous and unmissable.\ The absence of any deviation is a genuine
constraint.

\subsection{Topological Charge: Confirmed Resolution}

The resolution rests on the nature of the electron's charge in the Aether
model.\ The revised Paper~4 \cite{otto-em} identifies charge as the
relative rotation coupling between the bound core ring and its surrounding
vortex---co-rotating for the electron ($q = -e$), counter-rotating for the
positron ($q = +e$).\ This coupling is purely topological: a discrete
relative winding sign with no continuous deformation between the two
configurations and no spatial distribution of charge density.

Charge is therefore a property of the topological defect, not of the
$386$~fm ring geometry.\ The ring carries mass (circulation KE) and spin
(angular momentum) distributed over its circumference, but the charge
itself has no spatial extent.\ The form-factor prediction is
$F(Q^2) = 1$ at all $Q$, consistent with electron-scattering measurements
up to $Q \sim 200$~GeV.

The distinction is analogous to that of a magnetic monopole: a monopole's
charge is topological (from the field configuration's winding number),
not geometric (from a spatial charge distribution).\ Its form factor is
$F(Q^2) = 1$ because there is no charge distribution to resolve---the
charge \emph{is} the topology.\ The same applies here: there is no
``charge distribution'' on the electron ring to resolve, because charge is
sourced by the relative winding of ring and vortex, not by anything
distributed around the ring itself.

Paper~4's description of the electric field as spherically symmetric is
therefore exact (not approximate): the symmetry arises from topology, and
the coherent spin-phase wave field radiated by the ring inherits
this exact spherical symmetry at the monopole level.

\subsection{Status}

The topological-versus-geometric question that was previously open is
answered by the revised Paper~4 \cite{otto-em}: charge is topological, the
form factor satisfies $F(Q^2) = 1$ exactly, and the vortex-ring model is
consistent with electron-scattering data at all tested momentum transfers.\
The quantitative Coulomb-constant derivation is developed in the following
section (Section~\ref{sec:coulomb-derivation}), where the same relative
winding that carries the charge acts as a quantized spin-phase pump; that
derivation is independent of---and does not affect---the topological
resolution of the form-factor challenge.


% ============================================================
\section{Derivation of Coulomb's Constant from the Ring's Spin-Phase Winding Pump}
\label{sec:coulomb-derivation}
% ============================================================

Paper~4~\cite{otto-em} develops the electric field as a phase-locked
interaction between two coherent sources of discrete topological sign,
deriving the Coulomb-force \emph{structure} ($1/r^2$ falloff, sign rule)
without committing to a specific source realization.\ This section
supplies the quantitative source from the bound core ring topology of
Sections~\ref{sec:formation}--\ref{sec:form-factor} and derives $k$.\ An
earlier version of this section modeled the source as an axially
oscillating piston radiating density (sound) waves; that route is closed
by the channel-selection result below, and by Paper~7's identification of
photons as spin-wave excitations of the spinor condensate
\cite{otto-quantum}, which places the electromagnetic sector in the spin
channel from the outset.

\subsection{Channel Selection: The Source Cannot Be a Density Monopole}
\label{sec:channel-selection}

For any closed surface around the particle, the monopole (isotropic,
$1/r^2$-force-producing) component of a radiated field is fixed by
conservation: the net flux equals the rate of change of the corresponding
\emph{content} enclosed.\ For density (sound) waves the content is
displaced volume, and for a structure made of flow the available volume
is pinned by an identity: the Bernoulli deficit volume of a flow
structure equals its flow energy divided by the bulk modulus,
$V_{\mathrm{def}} = E_{\mathrm{flow}}/\Ka$.\ Matching the observed
Coulomb force through the acoustic (secondary Bjerknes) route
\cite{bjerknes1906} requires a volume amplitude
$V_0 = \sqrt{8\pi k e^2/(\rhoa\,\omega^2)} \approx 1.4\times10^{-40}$~m$^3$,
whose energy cost is
%
\begin{equation}
  \Ka V_0 \;\approx\; (75\text{--}110)\; m_e c^2
  \label{eq:density-nogo}
\end{equation}
%
---two orders of magnitude beyond the electron's entire energy budget,
which the mass equation of Section~\ref{sec:combined} has already spent
on the circulation.\ (Motion-based sources fare no better: a translating
or flexing structure that conserves its volume radiates as a dipole at
every instant, whatever its stiffness.)\ The density channel therefore
cannot host the Coulomb source.\ The spin channel can: the transverse
spin $S_+ = \psi_\uparrow^* \psi_\downarrow$ is a complex field whose
phase is \emph{compact}---a precessing source pumps it indefinitely
without accumulating content---so monopole sources of spin phase are not
content-bounded.\ The density no-go thus acts as a channel-selection
theorem, and it doubles as a structural reason why the density-channel
residual (gravity, Paper~3~\cite{otto-gravity}) is so much weaker than
electromagnetism.

\subsection{Spin-Phase Stiffness of the Condensate}
\label{sec:spin-stiffness}

Write the spinor condensate as $\Psi = (\psi_\uparrow, \psi_\downarrow)$
with $\psi_i = \sqrt{n_i}\,e^{i\theta_i}$ over a balanced background
($n_\uparrow = n_\downarrow = n/2$, $n = \rhoa/m_A$), and define the
total phase $\theta = (\theta_\uparrow + \theta_\downarrow)/2$ and the
relative (spin) phase $\varphi = \theta_\uparrow - \theta_\downarrow$.\
The Gross--Pitaevskii kinetic energy splits exactly:
%
\begin{equation}
  \frac{\hbar^2}{2 m_A}\!\left[n_\uparrow(\nabla\theta_\uparrow)^2
  + n_\downarrow(\nabla\theta_\downarrow)^2\right]
  = \frac{K_\theta}{2}(\nabla\theta)^2
  + \frac{K_\varphi}{2}(\nabla\varphi)^2,
  \qquad
  K_\varphi = \frac{\rhoa \hbar^2}{4 m_A^2}
            = \frac{\rhoa\,\kappa_e^2}{4\pi^2}.
  \label{eq:spin-stiffness}
\end{equation}
%
The spin-phase stiffness $K_\varphi$ is purely kinetic---the interaction
terms set wave speeds and healing lengths but do not enter
Eq.~(\ref{eq:spin-stiffness})---so the static results below carry no new
parameters, and the factor of~$4$ that previously had to be supplied by
the half-quantum assumption emerges here from the two-component
kinematics ($\theta_i = \theta \pm \varphi/2$).

\subsection{Flux Sources and the Interaction}
\label{sec:spin-force}

A particle that pumps relative phase is a flux source of the $\varphi$
field, with source strength
$q = \oint \nabla\varphi \cdot d\mathbf{A}$; by two-component
continuity, $q = (2 m_A/\hbar n)\,d(\Delta N)/dt$, where $\Delta N$ is
the enclosed spin-population imbalance.\ The energy functional
(\ref{eq:spin-stiffness}) is mathematically identical to incompressible
potential flow under $\rho \leftrightarrow K_\varphi$ and volume flux
$\leftrightarrow q$, so potential theory gives the interaction of two
sources directly: $E_{\mathrm{int}} = K_\varphi q_1 q_2/(4\pi d)$, and
for two synchronized oscillating pumps $q_i(t) = q_0 \sin(\pm\omega t)$
the time-averaged force is
%
\begin{equation}
  F \;=\; \frac{K_\varphi\, q_1 q_2}{8\pi\, d^{\,2}}.
  \label{eq:spin-force}
\end{equation}
%
The sign structure is that of flux sources of a gradient field (the
vortex-gas rule): like signs repel, opposite signs attract.\ The sign of
each pump is its winding sense---co-rotating versus counter-rotating
core, exactly the topological charge of
Section~\ref{sec:form-factor}---so Eq.~(\ref{eq:spin-force}) is
electrostatics with Gauss-law charges: field lines are steady
spin-current streamlines, and Gauss's law is pump conservation.

\subsection{The Winding Pump}
\label{sec:winding-pump}

The source amplitude follows from the ring's two rotations
(§\ref{sec:two-clocks}).\ The bound core co-rotates with the host flow
that created it, so the winding rotates poloidally---about the core
tube's axis---at the inherited rate
%
\begin{equation}
  \Omega \;=\; \frac{v(\xi)}{\xi} \;=\; \frac{m_A c^2}{\hbar},
  \label{eq:pump-clock}
\end{equation}
%
a property of the medium ($v(\xi) = c/\sqrt{2}$ and $\xi$ scale so that
$m_A$ cancels from the construction), and hence common to every
singly-wound structure---which is what allows unequal-mass charges to
remain synchronized.\ A rigidly rotating pattern changes no enclosed
totals; what pumps is the rotation's effect at the boundary: the $m = 1$
winding rotating at $\Omega$ makes the \emph{local} relative phase at
every boundary point precess uniformly at $\Omega$ (an axisymmetric,
monopole-compatible drive), which is an AC Josephson condition with bias
$\hbar\Omega = m_A c^2$---one constituent rest energy, the DCE
bind--unbind scale of Section~\ref{sec:threshold}.\ The boundary layer
therefore converts spin population sinusoidally at $\Omega$.\ The pump
amplitude is normalized by the same topological postulate that fixes the
effective current in §\ref{sec:bohr-magneton}: \emph{one winding quantum
per cycle}, giving
%
\begin{equation}
  q_0 \;=\; 2\pi\,\rloop,
  \label{eq:q0}
\end{equation}
%
the ring's own circumference.\ Three computed checks support this
normalization without yet deriving it: the required conversion current
is $I = n\,\kappa_e\,\rloop$ (frequency-independent); the boundary
layer's required duty is $2/\pi$ of its population per turn (order
unity---compare the factor $\sim\!300$ by which the density channel
falls short); and the layer then operates at $\approx 1/(2\pi)$ of its
critical spin current, i.e.\ well subcritical, which is what makes the
pumping sinusoidal.\ The remaining first-principles gap is the
boundary-layer conversion calculation itself
(Section~\ref{sec:open-questions}).

\subsection{The Coulomb Constant}
\label{sec:k-result}

Combining Eqs.~(\ref{eq:spin-stiffness}), (\ref{eq:spin-force}), and
(\ref{eq:q0}), the frequency, the amplitude conventions, and all
effective areas cancel, leaving
%
\begin{equation}
  \boxed{\;k\,e^2 \;=\; \frac{K_\varphi\,q_0^2}{8\pi}
  \;=\; \frac{\rhoa\,\kappa_e^2\,\rloop^2}{8\pi}\;}
  \label{eq:k-final}
\end{equation}
%
---Coulomb's constant from the medium density, the circulation quantum,
and the loop radius alone.\ Eliminating $\rhoa\kappa_e^2\rloop$ with the
mass equation (Section~\ref{sec:combined}) and using
$m_e c^2\, \rloop = \hbar c$ gives the dimensionless content of the
joint model:
%
\begin{equation}
  \alpha^{-1} \;=\; 4\pi \left[\ln\!\left(\frac{8\sqrt{2}\,m_A}{m_e}\right)
  - \tfrac{1}{2}\right]\left(1 + f_{\mathrm{core}} - 0.017\right),
  \label{eq:alpha-master}
\end{equation}
%
in which $\rhoa$ cancels entirely and $m_A$ enters only
logarithmically.\ Numerically, Eq.~(\ref{eq:k-final}) reproduces the
observed $k = 8.99\times10^{9}$~N\,m$^2$/C$^2$ to $0.2\%$ at
$\rhoa = 5\times10^{11}$~kg/m$^3$ (the reference density of
Section~\ref{sec:ke-only}) and $m_A = 666$~MeV/$c^2$.\ The substantive
tests are not this evaluation, which selects the point, but the joint
consistency of Section~\ref{sec:k-cross-check} and the parameter-free
magneton match of §\ref{sec:bohr-magneton}, which shares the
normalization postulate of Eq.~(\ref{eq:q0}).

\subsection{Joint Consistency with the Electron Mass}
\label{sec:k-cross-check}

Both the Coulomb constant and the electron mass pin the same combination
of medium parameters.\ Eq.~(\ref{eq:k-final}) fixes
%
\begin{equation}
  K_C \;\equiv\; \frac{\rhoa}{m_A^2}
  \;=\; \frac{8\, m_e^2 c^2 e^2 k}{\pi \hbar^4}
  \;\approx\; 3.54 \times 10^{65}\ \mathrm{kg^{-1}\,m^{-3}},
  \label{eq:k-C}
\end{equation}
%
while the mass equation fixes $\rhoa/m_A^2$ up to the core-energy
correction.\ Requiring both---equivalently, solving
Eq.~(\ref{eq:alpha-master}) for the observed $\alpha$---demands
%
\begin{equation}
  f_{\mathrm{core}} \;\approx\; 0.18\text{--}0.26
  \qquad (0.22\ \text{at}\ m_A = 666\ \mathrm{MeV}/c^2),
  \label{eq:fcore-requirement}
\end{equation}
%
somewhat above the a-priori estimate
$f_{\mathrm{core}} \sim 5$--$14\%$ of Section~\ref{sec:core-energy}.\
With that a-priori band the predicted $k$ runs $7$--$17\%$ high across
the convergence zone; exact agreement is a sharp, pre-registered
requirement on the GPE core-energy calculation.\ The joint consistency
is a \emph{requirement on} $f_{\mathrm{core}}$, not an independent
overdetermination---but it is now stated at a single common density, and
the exact-agreement point ($m_A \approx 666$~MeV/$c^2$ at
$\rhoa \approx 5.0\times10^{11}$~kg/m$^3$) coincides with the reference
density of the mass analysis rather than sitting at a different one.\
The requirement is also sensitive to the vortex-energy bracket constant
inherited from Paper~7 (the $-\tfrac12$ in Eq.~\ref{eq:paper7-mass});
verifying that constant against the vortex-ring literature shifts
Eq.~(\ref{eq:fcore-requirement}) accordingly and is flagged in
Section~\ref{sec:k-flags}.

\subsection{Parameter Constraints}
\label{sec:k-tightened}

Eq.~(\ref{eq:k-C}) collapses the two-dimensional $(\rhoa, m_A)$ space to
a curve.\ Along it:

\begin{center}
\begin{tabular}{@{}cccc@{}}
\toprule
$m_A$ (MeV/$c^2$) & $\rhoa$ (kg/m$^3$) & $\xi$ (fm) & $\Ka$ (J/m$^3$) \\
\midrule
$500$ & $2.8\times10^{11}$ & $0.28$ & $2.5\times10^{28}$ \\
$666$ & $5.0\times10^{11}$ & $0.21$ & $4.5\times10^{28}$ \\
$850$ & $8.1\times10^{11}$ & $0.16$ & $7.3\times10^{28}$ \\
\bottomrule
\end{tabular}
\end{center}

The combination
$\xi^2 \Ka = \hbar^2 K_C/2 \approx 1.97\times10^{-3}$~J/m is exactly
invariant along the curve.\ Evaluating the compression ratio of Paper~7
with the lattice-QCD string tension
$\sigma \approx 1.6\times10^{5}$~J/m gives
$X_{\mathrm{Coulomb}} = 8\sigma/(\pi\hbar^2 K_C) \approx 1.0\times10^{8}$
($\log_{10} X \approx 8.0$), about $50\%$ above Paper~7's
electron-mass-only estimate $X \approx 7\times10^{7}$, with bridge sound
speed $c_s = c/\sqrt{X} \approx 2.9\times10^{4}$~m/s; the shift does not
overturn any qualitative strong-force conclusion.\ Final pinning of the
point on the curve awaits the $f_{\mathrm{core}}$ calculation
(Eq.~\ref{eq:fcore-requirement}).

\subsection{Honest Flags}
\label{sec:k-flags}

The derivation rests on stated, checkable assumptions:

\begin{enumerate}[leftmargin=2em]
  \item \textbf{Pump normalization (one winding per cycle).}\
    Eq.~(\ref{eq:q0}) is the same postulate that yields the Bohr
    magneton exactly (§\ref{sec:bohr-magneton}), and the capacity and
    operating-point checks of §\ref{sec:winding-pump} pass, but the
    boundary-layer conversion (duty $2/\pi$ per turn) has not been
    derived from the GPE.
  \item \textbf{Species inheritance of the charge quantum.}\
    Eq.~(\ref{eq:q0}) uses the \emph{electron's} loop radius.\ Because
    $m_e$ is itself fixed by the medium (Section~\ref{sec:combined}),
    $q_0 = 2\pi\hbar/(m_e c)$ can be read as a medium constant---the
    charge quantum of the condensate's ground ring mode---but the
    mechanism by which a muon or proton inherits exactly this quantum
    is open (§\ref{sec:species-scaling}).
  \item \textbf{Core-energy requirement.}\
    Eq.~(\ref{eq:fcore-requirement}) demands
    $f_{\mathrm{core}} \approx 0.18$--$0.26$, above the rough
    $1/\ln(\rloop/\xi)$ estimate of Section~\ref{sec:core-energy}; the
    GPE calculation decides.\ The target moves if the vortex-energy
    bracket constant ($-\tfrac12$) is revised against the literature.
  \item \textbf{Synchronization.}\ Universal phase-locking of the pumps
    (in-phase for like signs) is taken from the injection-locking
    behavior of mutually coupled self-sustained oscillators and is not
    derived here.
  \item \textbf{Isolated-charge energetics.}\ The reactive (standing)
    character of a single charge's spin-phase field, as opposed to
    radiative loss at $\Omega$, requires an explicit treatment.
  \item \textbf{Spin-sector mode bookkeeping.}\ The statics above use
    only the kinetic stiffness (\ref{eq:spin-stiffness}), but
    photons-at-$c$ requires the magnon branch at $c$ while the
    cosmological spin-nematic modes of Paper~6 are slow; the mode
    inventory should be made explicit.
\end{enumerate}


% ============================================================
\section{Structural Confirmation: The Bohr Magneton from Ring Geometry}
\label{sec:bohr-magneton}
% ============================================================

The derivation of §\ref{sec:coulomb-derivation} fixes the Coulomb
constant $k$ from the bound core ring's spin-phase source amplitude.\ A
complementary, parameter-free check on the ring geometry comes from
computing its classical magnetic moment and comparing to the observed
Bohr magneton $\mu_B = e\hbar/(2m_e)$.\ The match is exact: no fit, no
closure assumptions, no source-identification factor.

\subsection{Effective Current of the Bound Core Ring}

The ring carries charge $e$ as a topological winding number
(§\ref{sec:form-factor}).\ The winding \emph{pattern} advances
azimuthally around the loop once per Compton period
$T_C = 2\pi\hbar/(m_e c^2)$.\ This advance is a phase motion, not a
material flow: the pattern speed at the loop radius is
$\omega_C\,\rloop = c$ exactly, while the host circulation at that
radius moves constituents at only
$v(\rloop) = m_e c/(2 m_A) \approx 4\times10^{-4}\,c$.\ (Were the
advance material transport, the combined azimuthal and poloidal speeds
of the core constituents would exceed $c$; the winding advance is the
motion of a pattern in the relative phase, not of the constituents that
carry it.)\ An external observer averaging over the loop therefore sees
one quantum of charge pass any azimuthal cross-section per Compton
period:
\begin{equation}
I_{\text{eff}} = \frac{e}{T_C}
              = \frac{e\,\omega_C}{2\pi}
              = \frac{e\,m_e c^2}{2\pi\,\hbar}.
\label{eq:I-eff-bohr}
\end{equation}

\subsection{Two Windings, Two Clocks}
\label{sec:two-clocks}

The azimuthal advance above is one of \emph{two} rotations the bound
structure carries, and they must not be conflated.\ The same
relative-phase winding also rotates poloidally---about the core tube's
own axis---at the rate inherited from the host vortex flow at the
formation boundary, $\Omega = v(\xi)/\xi = m_A c^2/\hbar$
(Eq.~\ref{eq:pump-clock}): a property of the medium, common to every
singly-wound structure regardless of its loop radius.\ The two clocks
differ by $\Omega/\omega_C = m_A/m_e \approx 1.3\times10^{3}$: the
winding re-threads poloidally some thirteen hundred times per azimuthal
cycle---a fine-pitch screw.\ The division of labor is clean.\ The
poloidal rotation is material (the bound core co-rotates with the flow
that created it, at $v(\xi) = c/\sqrt{2}$, safely subluminal) and is the
clock of the Coulomb source (§\ref{sec:winding-pump}).\ The azimuthal
advance is pattern motion and is the clock of the magnetic moment
computed here.\ The poloidal rotation is also compatible with everything
below: its local rotation axes point azimuthally and close around the
loop, so it contributes no net angular momentum about the symmetry axis
and leaves both $\mu_M$ and $S$ untouched.

\subsection{Magnetic Moment of the Loop}

Treating the ring as a closed current loop of radius $\rloop$ and area
$\mathcal{A} = \pi \rloop^2$, the classical magnetic moment is
$\mu_M = I_{\text{eff}}\,\mathcal{A}$.\ Substituting
$\rloop = \hbar/(m_e c)$ from Paper~7:
\begin{equation}
\mu_M = \frac{e\,m_e c^2}{2\pi\,\hbar}\cdot \pi\cdot\frac{\hbar^2}{m_e^2 c^2}
      = \frac{e\,\hbar}{2 m_e}
      = \mu_B.
\label{eq:bohr-magneton}
\end{equation}
The cancellation in~(\ref{eq:bohr-magneton}) is exact and
parameter-free.

\subsection{What the Match Confirms}

Equation~(\ref{eq:bohr-magneton}) requires three
independently-motivated structural inputs to combine in precisely the
right way:
\begin{itemize}
\item The ring radius is the reduced Compton wavelength:
  $\rloop = \hbar/(m_e c)$.\ Any other length scale produces the wrong
  answer.
\item The azimuthal clock is the Compton frequency:
  $\omega_C = m_e c^2/\hbar$.\ This is the frequency of the winding
  pattern's advance around the loop (Paper~7), distinct from the
  poloidal clock $\Omega$ of §\ref{sec:two-clocks}.
\item The effective current is one charge per period:
  $I_{\text{eff}} = e\,\omega_C/(2\pi)$.\ This follows from the
  topological character of charge established in
  §\ref{sec:form-factor}: the ring carries a single winding whose
  azimuthal pattern advances exactly once per cycle.
\end{itemize}
Each input is motivated independently elsewhere in the framework, and
their product yielding exactly $\mu_B$ is a non-trivial geometric
self-consistency.\ As developed below, the third input is not unique to
this calculation: the same one-winding-per-cycle normalization
independently fixes the Coulomb source amplitude, making the pair of
results a single-postulate structure rather than a coincidence.

\subsection{Gyromagnetic Ratio $g = 2$}

The same construction gives the Dirac value of the gyromagnetic ratio
without further assumptions.\ The bound core ring carries quantized
half-circulation $\kappa_e = h/(2 m_A)$
(cf.\ Paper~7~\cite{otto-quantum}), which corresponds to spin angular
momentum $S = \hbar/2$.\ Combining with~(\ref{eq:bohr-magneton}):
\begin{equation}
\mu_B \;=\; g\cdot\frac{e}{2 m_e}\cdot S \quad\Longrightarrow\quad
g \;=\; 2,
\end{equation}
to leading order.\ (The identification is exact per constituent---each
participant in the half-quantum flow carries
$m_A \kappa_e/2\pi = \hbar/2$---while the bookkeeping that assigns the
\emph{structure} a total of exactly $\hbar/2$ is carried by Paper~7's
quantized-circulation argument and is not re-derived here.)\ The
standard QED anomalous corrections
$g = 2 + \alpha/\pi + \mathcal{O}(\alpha^2)$ are expected to arise from
coupled fluctuations of the surrounding Aether condensate (the analog
of the radiative diagrams in QED) and lie beyond the classical ring
picture.

\subsection{What This Does \emph{Not} Determine}

The match in~(\ref{eq:bohr-magneton}) yields $\mu_B$ in SI units of
$\mathrm{A\cdot m^2}$ from purely geometric and kinematic ring
properties.\ It does \emph{not} fix the vacuum permeability $\mu_0$.\
That constant enters only when one computes the magnetic field produced
by the moment at distance $r$ ($B = \mu_0\mu_M/(4\pi r^3)$) or the
interaction energy of two such moments, and it is fixed
($\mu_0 = 4\pi k/c^2$) by the Coulomb-constant derivation of
§\ref{sec:coulomb-derivation}, whose source amplitude is normalized by
the \emph{same} postulate used in Eq.~(\ref{eq:I-eff-bohr}): one winding
advance per Compton cycle, giving pump strength
$q_0 = 2\pi\,\rloop$---one winding quantum across the ring aperture per
cycle.\ The two calculations therefore share their single dimensionless
input:
\begin{itemize}
\item one winding per cycle, counted as an azimuthal current
  $\;\rightarrow\;$ $\mu_B$, exactly;
\item one winding per cycle, counted as a spin-phase pump
  $\;\rightarrow\;$ $k$, to $0.2\%$ at the reference point of
  §\ref{sec:k-result}.
\end{itemize}
One topological postulate, two observables.\ This is the sense in which
the product in the preceding subsection is quantized rather than tuned:
a normalization adjusted by hand to fit $k$ would have no reason to
land the magneton exactly, and vice versa.

\subsection{Species Scaling: Why $\mu$ Scales and $e$ Must Not}
\label{sec:species-scaling}

The construction extends across the charged leptons by
$\rloop \to \hbar/(m c)$ and $T_C \to 2\pi\hbar/(m c^2)$, predicting
$\mu = e\hbar/(2m)$---as observed: the muon's moment is
$e\hbar/(2m_\mu)$ with the same $g \approx 2$.\ The scaling succeeds
precisely because $e$ enters~(\ref{eq:I-eff-bohr}) as a universal input
while only the geometry scales with the species mass.\ This is a
structural constraint on the companion Coulomb derivation: the charge
quantum must \emph{not} inherit the $\rloop$ scaling that the moment
correctly does---a source strength proportional to each species' own
loop radius would make the muon's charge smaller than the electron's by
$m_e/m_\mu$, against observation.\ The resolution direction is
available within the framework: $m_e$ is itself fixed by the medium
through the self-consistent mass equation (§\ref{sec:combined}), so
$q_0 = 2\pi\hbar/(m_e c)$ is a constant of the medium---the charge
quantum of the condensate's ground ring mode---to be inherited by any
charged structure, whatever its own size and clocks.\ Making that
inheritance mechanism explicit is an open problem
(§\ref{sec:open-questions}).

\subsection{Status}

\textbf{Derived (D).}\ Given (i) the ring radius from Paper~7, (ii) the
topological character of charge (§\ref{sec:form-factor}), and (iii) the
half-quantum circulation from Paper~7, the Bohr magneton and $g = 2$
follow with no free parameters or fits.\ The shared
one-winding-per-cycle normalization ties this result to the Coulomb
source amplitude, elevating the two to a single-postulate pair.\
\textbf{Open:} the total angular-momentum bookkeeping behind
$S = \hbar/2$, and the species-inheritance mechanism for the charge
quantum (§\ref{sec:species-scaling}).


% ============================================================
\section{Open Questions}
\label{sec:open-questions}
% ============================================================

\begin{enumerate}[leftmargin=2em]
  \item \textbf{DCE permanence threshold.}\ What is $v_{\mathrm{perm}}$?\ A
    first-principles DCE calculation in vortex geometry would determine the
    equilibrium core size and the KE/ring mass partitioning.\ This is the
    key open calculation needed to make the model fully quantitative.

  \item \textbf{Core energy in the logarithmic potential.}\ Computing the
    gradient and potential energy of a vortex core in the logarithmic
    condensate would determine $f_{\mathrm{core}}$ and, combined with the
    mass equation, pin down $m_A$ uniquely.\ The Coulomb constraint makes
    this fully quantitative: joint consistency requires
    $f_{\mathrm{core}} \approx 0.18$--$0.26$
    (Eq.~\ref{eq:fcore-requirement}), a pre-registered target for the
    calculation.

  \item \textbf{Muon and tau.}\ Could higher-energy vortex configurations
    (multiply-wound cores, knotted rings) stabilize at different masses
    through the same self-regulation mechanism?\ This is
    Paper~7~\cite{otto-quantum}, Open Question~11.

  \item \textbf{Form factor---residual geometric corrections.}\ See
    Section~\ref{sec:form-factor} for the full analysis.\ The revised
    Paper~4 \cite{otto-em} identifies charge as the topological rotation
    coupling between ring and vortex, which yields $F(Q^2) = 1$ exactly at
    the monopole level.\ A narrower remaining question is whether
    subleading geometric corrections of order $(Q r_{\mathrm{loop}})^2$ arise
    from the ring's finite extent in higher multipole moments of the
    coherent spin-phase wave field.

  \item \textbf{Winding-pump closure.}\
    Section~\ref{sec:coulomb-derivation} derives the Coulomb constant
    from the spin-phase winding pump with one normalization postulate
    ($q_0 = 2\pi\rloop$; one winding quantum per cycle,
    Eq.~\ref{eq:q0}).\ Remaining first-principles work:
    (a) the boundary-layer conversion calculation---the driven
    two-component GP problem at the rotating winding boundary, with
    pre-registered answer: duty $2/\pi$ of the layer population per
    poloidal turn; (b) the species-inheritance mechanism for the charge
    quantum (§\ref{sec:species-scaling}); (c) universal synchronization
    of the pumps from injection locking; (d) the isolated-charge
    energetics of the spin-phase field.

  \item \textbf{Cosmological formation.}\ What conditions in the early
    universe create vortex rings with the minimum circulation quantum?\ What
    determines the initial $\rloop$, and how does it relax to
    $\hbar/(m_e c)$?
\end{enumerate}


% ============================================================
\begin{thebibliography}{99}

\bibitem{otto-overview}
E.~Otto,
``Unifying Theory of Dynamic Aether Mechanics: General Overview,''
this series, Paper~1.

\bibitem{otto-relativity}
E.~Otto,
``Unifying Theory of Dynamic Aether Mechanics: Special Relativity,''
this series, Paper~2.

\bibitem{otto-gravity}
E.~Otto,
``Unifying Theory of Dynamic Aether Mechanics: Gravity,''
this series, Paper~3.

\bibitem{otto-em}
E.~Otto,
``Unifying Theory of Dynamic Aether Mechanics: Magnetic and Electric Fields,''
this series, Paper~4.

\bibitem{otto-strong}
E.~Otto,
``Unifying Theory of Dynamic Aether Mechanics: Strong and Weak Forces,''
this series, Paper~5.

\bibitem{otto-quantum}
E.~Otto,
``Unifying Theory of Dynamic Aether Mechanics: Quantum Mechanics,''
this series, Paper~7.

\bibitem{donnelly1991}
R.~J.~Donnelly,
\textit{Quantized Vortices in Helium~II}
(Cambridge University Press, 1991).

\bibitem{wilson2011}
C.~M.~Wilson \textit{et al.},
``Observation of the dynamical Casimir effect in a superconducting circuit,''
\textit{Nature} \textbf{479}, 376 (2011).

\bibitem{LEP-Bhabha}
OPAL Collaboration,
``Measurement of the running of the QED coupling in small-angle Bhabha
scattering at LEP,''
\textit{Eur.\ Phys.\ J.\ C} \textbf{45}, 1 (2006).

\bibitem{bjerknes1906}
V.~F.~K.~Bjerknes,
\textit{Fields of Force}
(Columbia University Press, New York, 1906).

\end{thebibliography}

\end{document}
